% NeurIPS Paper Skeleton
% Replace placeholders marked with <...> before use
\documentclass{article}

% NeurIPS style — update year as needed
\usepackage[final]{neurips_2025}

\usepackage[utf8]{inputenc}
\usepackage[T1]{fontenc}
\usepackage{hyperref}
\usepackage{url}
\usepackage{booktabs}
\usepackage{amsfonts}
\usepackage{amsmath}
\usepackage{amssymb}
\usepackage{nicefrac}
\usepackage{microtype}
\usepackage{graphicx}
\usepackage{xcolor}

% Custom commands
\newcommand{\ours}{\textsc{<MethodName>}}
\newcommand{\todo}[1]{\textcolor{red}{[TODO: #1]}}

\title{<Paper Title>}

\author{
  <Author 1> \\
  <Affiliation 1> \\
  \texttt{<email1>} \\
  \And
  <Author 2> \\
  <Affiliation 2> \\
  \texttt{<email2>} \\
}

\begin{document}

\maketitle

\begin{abstract}
% 4-sentence structure:
% S1: Problem context and importance
% S2: Key limitation of existing approaches
% S3: Our approach and key insight
% S4: Main results and significance
\todo{Write abstract — 150-250 words}
\end{abstract}

%% ============================================================
\section{Introduction}
\label{sec:intro}

% P1: Broad context — why this area matters
\todo{Broad motivation}

% P2: Specific problem — what exactly are we solving
\todo{Specific problem statement}

% P3: Existing approaches and their limitations
\todo{Prior work limitations}

% P4: Our approach — key insight and method overview
\todo{Our approach overview}

% P5: Contributions list
Our main contributions are:
\begin{itemize}
    \item \todo{Contribution 1 — conceptual/methodological}
    \item \todo{Contribution 2 — technical}
    \item \todo{Contribution 3 — empirical}
\end{itemize}

%% ============================================================
\section{Related Work}
\label{sec:related}

% Group by topic, not chronologically
\paragraph{<Topic 1>.}
\todo{Related work group 1}

\paragraph{<Topic 2>.}
\todo{Related work group 2}

\paragraph{<Topic 3>.}
\todo{Related work group 3}

%% ============================================================
\section{Method}
\label{sec:method}

% Overview paragraph — the "30-second version"
\todo{Method overview — reference Figure 1}

\subsection{Problem Formulation}
\label{sec:formulation}
\todo{Formal problem definition, notation}

\subsection{<Component 1>}
\label{sec:component1}
\todo{First key component}

\subsection{<Component 2>}
\label{sec:component2}
\todo{Second key component}

\subsection{Training and Inference}
\label{sec:training}
\todo{Training procedure, loss functions, inference details}

%% ============================================================
\section{Experiments}
\label{sec:experiments}

\subsection{Experimental Setup}
\label{sec:setup}

\paragraph{Datasets.}
\todo{Dataset descriptions and statistics}

\paragraph{Baselines.}
\todo{Baseline methods and their configurations}

\paragraph{Implementation details.}
\todo{Hyperparameters, hardware, training details}

\subsection{Main Results}
\label{sec:main_results}

\todo{Main comparison table and analysis}

% Example table structure:
% \begin{table}[t]
% \caption{Main results on <benchmark>. Best in \textbf{bold}, second best \underline{underlined}.}
% \label{tab:main}
% \centering
% \begin{tabular}{lcc}
% \toprule
% Method & Metric 1 & Metric 2 \\
% \midrule
% Baseline 1 & 00.0 & 00.0 \\
% Baseline 2 & 00.0 & 00.0 \\
% \midrule
% \ours & \textbf{00.0} & \textbf{00.0} \\
% \bottomrule
% \end{tabular}
% \end{table}

\subsection{Ablation Study}
\label{sec:ablation}
\todo{Ablation experiments — one table per research question}

\subsection{Analysis}
\label{sec:analysis}
\todo{Qualitative analysis, visualizations, failure cases}

%% ============================================================
\section{Conclusion}
\label{sec:conclusion}

% P1: Summarize what we did and key findings
\todo{Summary}

% P2: Limitations
\todo{Limitations}

% P3: Future work
\todo{Future directions}

%% ============================================================
% Acknowledgments (only in camera-ready)
% \begin{ack}
% \todo{Acknowledgments}
% \end{ack}

\bibliographystyle{plainnat}
\bibliography{references}

%% ============================================================
% Appendix
\appendix
\section{Additional Experimental Details}
\label{app:details}
\todo{Full hyperparameter tables, additional results}

\section{Proofs}
\label{app:proofs}
\todo{Deferred proofs}

\end{document}
